% Options for packages loaded elsewhere
\PassOptionsToPackage{unicode}{hyperref}
\PassOptionsToPackage{hyphens}{url}
\documentclass[
]{article}
\usepackage{xcolor}
\usepackage{amsmath,amssymb}
\setcounter{secnumdepth}{-\maxdimen} % remove section numbering
\usepackage{iftex}
\ifPDFTeX
  \usepackage[T1]{fontenc}
  \usepackage[utf8]{inputenc}
  \usepackage{textcomp} % provide euro and other symbols
\else % if luatex or xetex
  \usepackage{unicode-math} % this also loads fontspec
  \defaultfontfeatures{Scale=MatchLowercase}
  \defaultfontfeatures[\rmfamily]{Ligatures=TeX,Scale=1}
\fi
\usepackage{lmodern}
\ifPDFTeX\else
  % xetex/luatex font selection
\fi
% Use upquote if available, for straight quotes in verbatim environments
\IfFileExists{upquote.sty}{\usepackage{upquote}}{}
\IfFileExists{microtype.sty}{% use microtype if available
  \usepackage[]{microtype}
  \UseMicrotypeSet[protrusion]{basicmath} % disable protrusion for tt fonts
}{}
\makeatletter
\@ifundefined{KOMAClassName}{% if non-KOMA class
  \IfFileExists{parskip.sty}{%
    \usepackage{parskip}
  }{% else
    \setlength{\parindent}{0pt}
    \setlength{\parskip}{6pt plus 2pt minus 1pt}}
}{% if KOMA class
  \KOMAoptions{parskip=half}}
\makeatother
\usepackage{longtable,booktabs,array}
\usepackage{caption}
\captionsetup[table]{skip=6pt}
\newcounter{none} % for unnumbered tables
\usepackage{calc} % for calculating minipage widths
% Correct order of tables after \paragraph or \subparagraph
\usepackage{etoolbox}
\makeatletter
\patchcmd\longtable{\par}{\if@noskipsec\mbox{}\fi\par}{}{}
\makeatother
% Allow footnotes in longtable head/foot
\IfFileExists{footnotehyper.sty}{\usepackage{footnotehyper}}{\usepackage{footnote}}
\makesavenoteenv{longtable}
\usepackage{graphicx}
\makeatletter
\newsavebox\pandoc@box
\newcommand*\pandocbounded[1]{% scales image to fit in text height/width
  \sbox\pandoc@box{#1}%
  \Gscale@div\@tempa{\textheight}{\dimexpr\ht\pandoc@box+\dp\pandoc@box\relax}%
  \Gscale@div\@tempb{\linewidth}{\wd\pandoc@box}%
  \ifdim\@tempb\p@<\@tempa\p@\let\@tempa\@tempb\fi% select the smaller of both
  \ifdim\@tempa\p@<\p@\scalebox{\@tempa}{\usebox\pandoc@box}%
  \else\usebox{\pandoc@box}%
  \fi%
}
% Set default figure placement to htbp
\def\fps@figure{htbp}
\makeatother
% definitions for citeproc citations
\NewDocumentCommand\citeproctext{}{}
\NewDocumentCommand\citeproc{mm}{%
  \begingroup\def\citeproctext{#2}\cite{#1}\endgroup}
\makeatletter
 % allow citations to break across lines
 \let\@cite@ofmt\@firstofone
 % avoid brackets around text for \cite:
 \def\@biblabel#1{}
 \def\@cite#1#2{{#1\if@tempswa , #2\fi}}
\makeatother
\newlength{\cslhangindent}
\setlength{\cslhangindent}{1.5em}
\newlength{\csllabelwidth}
\setlength{\csllabelwidth}{3em}
\newenvironment{CSLReferences}[2] % #1 hanging-indent, #2 entry-spacing
 {\begin{list}{}{%
  \setlength{\itemindent}{0pt}
  \setlength{\leftmargin}{0pt}
  \setlength{\parsep}{0pt}
  % turn on hanging indent if param 1 is 1
  \ifodd #1
   \setlength{\leftmargin}{\cslhangindent}
   \setlength{\itemindent}{-1\cslhangindent}
  \fi
  % set entry spacing
  \setlength{\itemsep}{#2\baselineskip}}}
 {\end{list}}
\usepackage{calc}
\newcommand{\CSLBlock}[1]{\hfill\break\parbox[t]{\linewidth}{\strut\ignorespaces#1\strut}}
\newcommand{\CSLLeftMargin}[1]{\parbox[t]{\csllabelwidth}{\strut#1\strut}}
\newcommand{\CSLRightInline}[1]{\parbox[t]{\linewidth - \csllabelwidth}{\strut#1\strut}}
\newcommand{\CSLIndent}[1]{\hspace{\cslhangindent}#1}
\setlength{\emergencystretch}{3em} % prevent overfull lines
\providecommand{\tightlist}{%
  \setlength{\itemsep}{0pt}\setlength{\parskip}{0pt}}
\makeatletter
\@ifpackageloaded{subcaption}{}{\usepackage{subcaption}}
\@ifpackageloaded{caption}{}{\usepackage{caption}}
\captionsetup[subfigure]{margin=0.5em}
\AtBeginDocument{%
\renewcommand*\figurename{Figure}
\renewcommand*\tablename{Table}
}
\AtBeginDocument{%
\renewcommand*\listfigurename{List of Figures}
\renewcommand*\listtablename{List of Tables}
}
\newcounter{pandoccrossref@subfigures@footnote@counter}
\newenvironment{pandoccrossrefsubfigures}{%
\setcounter{pandoccrossref@subfigures@footnote@counter}{0}
\begin{figure}\centering%
\gdef\global@pandoccrossref@subfigures@footnotes{}%
\DeclareRobustCommand{\footnote}[1]{\footnotemark%
\stepcounter{pandoccrossref@subfigures@footnote@counter}%
\ifx\global@pandoccrossref@subfigures@footnotes\empty%
\gdef\global@pandoccrossref@subfigures@footnotes{{##1}}%
\else%
\g@addto@macro\global@pandoccrossref@subfigures@footnotes{, {##1}}%
\fi}}%
{\end{figure}%
\addtocounter{footnote}{-\value{pandoccrossref@subfigures@footnote@counter}}
\@for\f:=\global@pandoccrossref@subfigures@footnotes\do{\stepcounter{footnote}\footnotetext{\f}}%
\gdef\global@pandoccrossref@subfigures@footnotes{}}
\@ifpackageloaded{float}{}{\usepackage{float}}
\floatstyle{ruled}
\@ifundefined{c@chapter}{\newfloat{codelisting}{h}{lop}}{\newfloat{codelisting}{h}{lop}[chapter]}
\floatname{codelisting}{Listing}
\newcommand*\listoflistings{\listof{codelisting}{List of Listings}}
\makeatother
\usepackage{bookmark}
\IfFileExists{xurl.sty}{\usepackage{xurl}}{} % add URL line breaks if available
\urlstyle{same}
% fallback for those not using the hyperref driver hyperxmp:
\makeatletter
\@ifundefined{xmpquote}{\newcommand{\xmpquote}[1]{#1}}{}
\makeatother
\hypersetup{
  pdftitle={Noise bends a polynomial fit},
  pdfauthor={Ada Lovelace (Analytical Engine Institute)},
  hidelinks,
  pdfcreator={LaTeX via pandoc}}

\title{Noise bends a polynomial fit}
\author{Ada Lovelace (Analytical Engine Institute)}
\date{}

\begin{document}
\maketitle
\begin{abstract}
We measure how noise bends a polynomial fit, and find that its error
grows as the square of the noise.
\end{abstract}

\section{Introduction}\label{introduction}

\subsection{State of the art}\label{introduction-state-of-the-art}

Polynomial fits are
standard\textsuperscript{\citeproc{ref-smith-2020-fits}{1}--\citeproc{ref-kim-2017-survey}{3}}.
Noise is known to bias
them\textsuperscript{\citeproc{ref-lee-2018-noise}{4}}. An earlier draft
said so. A lost work claimed it first {[}to be resolved{]}, and a late
one agreed\textsuperscript{\citeproc{ref-late-work}{5}}. A third said it
again.

\section{Results}\label{results}

\subsection{How does noise change the fit?}\label{O1}

\subsubsection{The error grows with the noise}\label{R1}

The error of the fit follows

\begin{equation}\protect\phantomsection\label{eq:error}{ e(\sigma) = 0.8\,\sigma^2 + 0.02 }\end{equation}

and Eq.~\ref{eq:error} holds at every noise level we ran, as
Fig.~\ref{fig:F1} shows. Table~\ref{tbl:T1} lists the errors. The fit
follows Smith's
method\textsuperscript{\citeproc{ref-smith-2020-fits}{1}}.

\begin{longtable}[]{@{}ll@{}}
\caption{The error of the fit at each noise
level.}\label{tbl:T1}\tabularnewline
\toprule\noalign{}
noise & rmse \\
\midrule\noalign{}
\endfirsthead
\toprule\noalign{}
noise & rmse \\
\midrule\noalign{}
\endhead
\bottomrule\noalign{}
\endlastfoot
0.1 & 0.03 \\
0.2 & 0.05 \\
0.4 & 0.15 \\
\end{longtable}

\begin{figure}
\centering
\pandocbounded{\includegraphics[keepaspectratio,alt={The error of the fit grows with the noise.}]{figures/F1.png}}
\caption{The error of the fit grows with the noise.}\label{fig:F1}
\end{figure}

\section{Discussion}\label{discussion}

\subsection{Limits}\label{discussion-limits}

The noise model is Gaussian, as Lee
assumed\textsuperscript{\citeproc{ref-lee-2018-noise}{4}}. A heavier
tail would bend the fit further, and the next runs should try one.

\section{Methods}\label{methods}

\subsection{Fitting}\label{methods-fitting}

We fit by least squares, as Jones
did\textsuperscript{\citeproc{ref-jones-2019-least}{2}}.

\section{Acknowledgments}\label{acknowledgments}

{[}TO BE GIVEN: acknowledgments{]}

\section{Funding}\label{funding}

{[}TO BE GIVEN: funding{]}

\section{Competing interests}\label{competing-interests}

The authors declare none.

\section{Data availability}\label{data-availability}

The data are synthetic and the project's runs made them.

\section{Code availability}\label{code-availability}

{[}TO BE GIVEN: the code repository, which setup has not settled{]}

\section{References}\label{references}

\protect\phantomsection\label{refs}
\begin{CSLReferences}{0}{0}
\bibitem[\citeproctext]{ref-smith-2020-fits}
\CSLLeftMargin{1. }%
\CSLRightInline{Smith, A. \& Li, B. Polynomial fits in practice.
\emph{Journal of Fits} \url{https://doi.org/10.5555/smith} (2020)
doi:\href{https://doi.org/10.5555/smith}{10.5555/smith}.}

\bibitem[\citeproctext]{ref-jones-2019-least}
\CSLLeftMargin{2. }%
\CSLRightInline{Jones, C. Least squares for curves. \emph{Numerical
Notes} \url{https://doi.org/10.5555/jones} (2019)
doi:\href{https://doi.org/10.5555/jones}{10.5555/jones}.}

\bibitem[\citeproctext]{ref-kim-2017-survey}
\CSLLeftMargin{3. }%
\CSLRightInline{Kim, F. A survey of fitting. \emph{Fitting Reviews}
\url{https://doi.org/10.5555/kim} (2017)
doi:\href{https://doi.org/10.5555/kim}{10.5555/kim}.}

\bibitem[\citeproctext]{ref-lee-2018-noise}
\CSLLeftMargin{4. }%
\CSLRightInline{Lee, D. \& Hall, E. Noise and bias in fits. \emph{arXiv}
Preprint at \url{https://arxiv.org/abs/1801.00001} (2018).}

\bibitem[\citeproctext]{ref-late-work}
\CSLLeftMargin{5. }%
\CSLRightInline{Late, A. A late agreement. \emph{Journal of Tests}
\url{https://doi.org/10.5555/late} (2024)
doi:\href{https://doi.org/10.5555/late}{10.5555/late}.}

\end{CSLReferences}

\end{document}
